% Sezione completa: sweep eseguito (job 34137, gnode01, 2026-09-09), dati in
% mandelbrot/data/run_34137.csv e mandelbrot/data/merged.csv.
% I limiti p/lambda_p provengono da mandelbrot/data/block_imbalance.csv.
\subsection{MPI: memoria distribuita}
\label{sec:mpi}

Il passaggio a MPI cambia il vincolo fondamentale, non la decomposizione. La
grana resta la riga, per le stesse ragioni della sezione~\ref{sec:openmp}, ma i
processi non condividono più uno spazio di indirizzamento: ciascun rank possiede
una propria copia parziale della matrice degli escape time, e il risultato
completo esiste solo dopo un atto esplicito di comunicazione. Ne discendono le
due differenze che governano l'intera sottosezione. Primo, la matrice va
\emph{assemblata}: nei due schemi statici con una \texttt{MPI\_Gatherv} finale
su rank $0$, nello schema dinamico incrementalmente, riga per riga, man mano che
i risultati arrivano. Secondo, e più importante, il ribilanciamento del carico
non è più gratuito: sotto OpenMP un thread inattivo prelevava lavoro da una coda
in memoria condivisa a costo trascurabile, mentre qui ogni atto di
ridistribuzione è un messaggio, potenzialmente attraverso la rete. È questo
prezzo a rendere il confronto fra schemi la vera domanda del paradigma.

Il parallelismo resta quello per righe, ma l'unità è ora il processo e non il
thread; il numero di unità di calcolo è spazzato sugli stessi valori
($2, 4, 8, 16, 32$) usati per OpenMP, e griglia e $N_{\max}$ restano fissi
($1024^2$, $1000$), così che il confronto fra paradigmi avvenga sullo stesso
carico di lavoro. Non si sfrutta la simmetria, coerentemente con tutti gli altri
benchmark di scaling: dimezzerebbe il lavoro senza cambiare la matrice
risultante, ma renderebbe metà delle righe una copia anziché un calcolo,
falsando il profilo $\{w_r\}$ su cui poggia l'intera analisi di sbilanciamento.

% ============================================================
\subsubsection{Design dello sweep}
\label{sec:mpi-sweep}
% ============================================================

La domanda progettuale non è quanto parallelizzare, ma \emph{con quale regola
assegnare le righe ai rank}. A differenza di OpenMP, dove la scelta era fra
politiche di scheduling dello stesso pool di thread, qui le regole differiscono
anche nel volume di comunicazione che generano: è possibile bilanciare meglio
pagando più messaggi, oppure bilanciare altrettanto bene senza aggiungerne.
Questo asse è ciò che i tre schemi esplorano.

\paragraph{\texttt{block}.}
Il rank $r$ possiede un intervallo contiguo di $H/p$ righe. La
distribuzione è nota a tutti prima di iniziare: non esiste dialogo, e la
comunicazione si riduce a un'unica \texttt{MPI\_Gatherv} finale. Poiché, per la
viewport scelta, il corpo del frattale occupa le righe centrali della griglia e
le righe di bordo contengono solo punti esterni che sfuggono in poche
iterazioni, i rank centrali concentrano quasi tutto il lavoro. La riga più
leggera è la riga $0$ ($2\,281$ iterazioni totali, circa $2$ per pixel) e la
più pesante la riga $511$, al centro ($700\,806$, circa $684$ per pixel): un
rapporto di $307\times$. I rank che ricevono il primo e l'ultimo blocco sono
necessariamente veloci per costruzione; i rank centrali portano da soli la
maggior parte del lavoro. È il caso di massimo sbilanciamento con comunicazione
minima, e il suo limite $S \le p/\lambda_p$ è calcolabile dal solo profilo
seriale prima ancora che il codice MPI esista.

\paragraph{\texttt{cyclic}.}
La riga $r$ va al rank $r \bmod p$. Ogni rank riceve righe sparse su tutta
l'immagine --- un po' di bordo, un po' di centro --- e il carico si media da sé,
perché il pattern di fuga varia su una scala spaziale molto più grande della
singola riga. L'assegnazione resta \emph{statica}: ogni rank sa già quali righe
possiede, non esiste coda residua, e la comunicazione è la stessa unica gather
del block --- a meno di un riordino locale su rank $0$ per de-interlacciare il
buffer, che non è un messaggio ma non è nemmeno gratuito, come mostreranno i
risultati.

Ciò che rende il cyclic interessante non è soltanto che batta il block: è che
\emph{esce dall'asse} bilanciamento--comunicazione su cui si trova il dinamico.
Lungo quell'asse, ridurre il chunk avvicina al bilanciamento ma moltiplica i
messaggi; il cyclic ottiene un bilanciamento quasi equivalente al dinamico con
zero messaggi aggiuntivi, perché è l'interlacciamento statico a mediare il
carico, non la coordinazione a runtime. La domanda aperta che l'esperimento deve
rispondere è se questo vantaggio sia sufficiente a battere o pareggiare il
master--worker nonostante il riordino.

\paragraph{\texttt{dynamic} (master--worker).}
Il rank $0$ non calcola nulla: tiene le righe non distribuite come coda di
lavoro, innesca ciascun worker con una riga e poi alimenta con la successiva il
primo worker che riporta un risultato. L'assegnazione si adatta al costo reale
di ogni riga --- chi è bloccato su una riga pesante semplicemente non viene
rifornito --- e il bilanciamento è perfetto per costruzione. Il prezzo è una
coppia di messaggi richiesta/risposta per ogni riga. Questo schema è il limite
opposto al block: bilanciamento massimo, overhead di comunicazione massimo.

È importante notare che, al crescere del chunk, il master--worker non degenera
nel cyclic ma nel block: con chunk pari a $H/p$ righe, dove $p$ è il numero di
worker, ogni worker riceve un unico
blocco contiguo, che è esattamente la decomposizione a blocchi contigui. Il
cyclic non è un punto sull'asse chunk--overhead: è un meccanismo ortogonale che
non occupa alcuna posizione su quell'asse.

\paragraph{Grana della concessione e conteggio delle unità.}
La concessione è fissata a una riga, senza sweep sul chunk, per tre ragioni.
Un chunk intermedio si colloca fra estremi già misurati: comunica più del
cyclic, che non aggiunge messaggi, e bilancia meno del chunk unitario. In secondo luogo l'estremo è il
punto più informativo: accoppia il miglior bilanciamento possibile con il massimo
costo di coordinamento, ed è precisamente il confronto che lo schema esiste per
esporre. Infine l'asse del chunk è già spazzato dove costa poco, nei run OpenMP
\texttt{dynamic,\{1,16,64\}}; tenendo il chunk unitario da entrambi i lati, i
due scheduler dinamici diventano la stessa politica su due modelli di memoria
diversi, e il loro divario isola il costo del passaggio di messaggi rispetto
all'incremento di un contatore condiviso.

Poiché il master non calcola, per lo schema dinamico $p$ indica il numero di
\emph{worker} e non di processi: gli sweep girano su $np \in \{3, 5, 9, 17,
33\}$ per mantenere le unità di calcolo a $\{2, 4, 8, 16, 32\}$, allineate agli
altri schemi e ai thread OpenMP. Va dichiarato che la contabilità è leggermente
generosa verso il master--worker: il cluster alloca $p+1$ core, uno dei quali
non calcola; speedup ed efficienza sono per unità di calcolo, e il costo reale
in risorse è $p+1$ processi.

% ============================================================
\subsubsection{Metriche specifiche}
\label{sec:mpi-metriche}
% ============================================================

Alle metriche comuni della sezione~\ref{sec:metriche-comuni} MPI aggiunge una
grandezza propria: la \emph{frazione di comunicazione}, cioè la quota di $T_\parallel(p)$
spesa a muovere dati anziché a calcolare. È la metrica che dà corpo al costo
introdotto dalla memoria distribuita, e la sua definizione operativa differisce
fra schemi statici e dinamico.

Negli schemi statici la comunicazione è un unico \texttt{MPI\_Gatherv} finale, e
il suo tempo è misurato \emph{dopo} una barriera. La scelta è deliberata: senza
barriera il rank in ritardo farebbe apparire dentro la gather l'attesa dovuta
allo sbilanciamento, gonfiando \texttt{comm\_fraction} con un tempo che non è
trasferimento di dati. Con la barriera l'inattività resta parcheggiata prima
dell'operazione --- e comunque dentro la regione cronometrata, quindi $T_\parallel(p)$
resta onesto --- mentre \texttt{comm\_fraction} misura solo movimento genuino di
dati. Lo sbilanciamento è riportato altrove, da $\lambda_p$ e dallo speedup; e
tenere le due cose separate è ciò che rende leggibile il confronto.

Nello schema dinamico la definizione è necessariamente diversa. Il master
comunica il $100\%$ del tempo per costruzione --- è la sua sola funzione ---
quindi il suo tempo non direbbe nulla sulla frazione di coordinamento pagata
\emph{dal calcolo}. Si misura invece, per ciascun worker, la frazione del proprio
tempo trascorsa dentro le chiamate MPI: l'invio del risultato di una riga e
l'attesa bloccante del comando successivo. Questa attesa è inclusa
deliberatamente, perché contiene sia la latenza del messaggio sia la contesa
quando più worker riportano insieme --- che è precisamente il prezzo della
coordinazione dinamica. Il valore medio dei worker è raccolto con una
\texttt{MPI\_Reduce} su rank $0$ e riportato come \texttt{comm\_fraction}. È
l'unico degli schemi per cui questa colonna diventa una metrica significativa:
negli statici la frazione è attesa prossima a zero e serve come controllo, nel
dinamico è il numero che quantifica il costo del self-scheduling.

% ============================================================
\subsubsection{Frammenti di codice rilevanti}
\label{sec:mpi-codice}
% ============================================================

Come per OpenMP, un solo binario copre tutti gli schemi: la variabile d'ambiente
\texttt{MPI\_DECOMP} seleziona a esecuzione fra \texttt{block}, \texttt{cyclic}
e \texttt{dynamic} (default \texttt{block}), e il job passa l'etichetta
corrispondente a \texttt{-{}-schedule} per il CSV. Il codice completo del kernel
si trova nel file \repofile{mandelbrot/src/mandelbrot_mpi.cpp}, mentre lo strato
driver descritto più avanti si trova nel file \repofile{mandelbrot/src/main.cpp}.

La differenza fra i due schemi statici è una sola espressione --- quale riga
globale corrisponda all'$i$-esima riga locale:

\begin{verbatim}
// block:  global row = row_displs[rank] + i   (contiguous range)
// cyclic: global row = rank + i*size          (strided)
const int row = (decomp == Decomp::Block) ? (row_displs[rank] + i)
                                          : (rank + i * size);
\end{verbatim}

È il frammento più significativo della sottosezione: il cyclic ottiene gran
parte del bilanciamento del master--worker cambiando un indice. I conteggi di
righe per rank sono identici nei due casi (\texttt{base} righe ciascuno, le
prime \texttt{rem} ne prendono una in più); cambia solo \emph{quali} righe un
rank possiede. Il block raccoglie direttamente nella matrice finale, essendo i
suoi dati già contigui; il cyclic raccoglie in un buffer compattato che il rank
$0$ poi de-interlaccia, un riordino locale in memoria che non aggiunge alcun
messaggio. La simmetria fra i due schemi è però meno perfetta di quanto la
singola espressione suggerisca: quel riordino è una copia seriale dell'intera
matrice, eseguita da un solo rank e di costo indipendente da $p$. Non entra in
\texttt{comm\_fraction} --- non è comunicazione --- ma entra in $T_\parallel(p)$,
ed è l'unico costo del cyclic assente negli altri due schemi; se ne discute
il peso nell'analisi dei risultati. Sarebbe eliminabile ricevendo direttamente
in posizione con un tipo derivato \texttt{MPI\_Type\_vector} di passo
\texttt{size}, al prezzo di una gather su tipo non contiguo.

La gather è cronometrata dopo una barriera, per la ragione discussa sopra:

\begin{verbatim}
MPI_Barrier(MPI_COMM_WORLD);      // park the straggler wait outside the timing
const double t0 = MPI_Wtime();    // of the gather (but inside T(p))
MPI_Gatherv(local.data(), my_rows * width, MPI_INT, recvbuf,
            px_counts.data(), px_displs.data(), MPI_INT, 0, MPI_COMM_WORLD);
g_last_comm_seconds = MPI_Wtime() - t0;
\end{verbatim}

Lo schema dinamico è invece l'unico che richiede un'architettura dedicata. Il
protocollo è minimale: il comando è un intero (l'indice di riga), il risultato è
\texttt{width+1} interi --- l'indice di riga seguito dagli escape time di quella
riga --- cosicché il master archivi il risultato all'arrivo senza tenere traccia
di chi detenga cosa. Il cuore è il ciclo di eventi del master:

\begin{verbatim}
while (active > 0) {
    MPI_Recv(result.data(), width + 1, MPI_INT, MPI_ANY_SOURCE, TAG_RESULT,
             MPI_COMM_WORLD, &status);        // whoever finishes first
    // file the row into the matrix, then feed that same worker
    if (next_row < height) {
        MPI_Send(&next_row, 1, MPI_INT, status.MPI_SOURCE, TAG_WORK, ...);
        ++next_row;
    } else {
        MPI_Send(&none, 1, MPI_INT, status.MPI_SOURCE, TAG_STOP, ...);
        --active;
    }
}
\end{verbatim}

\texttt{MPI\_ANY\_SOURCE} è ciò che realizza il self-scheduling: non è il master
a decidere chi lavora, ma il primo worker che si libera a farsi avanti. La
terminazione è implicita: ogni worker attivo detiene esattamente una riga in
sospeso, quindi \texttt{active} si annulla quando tutte le righe
sono state ricevute. Non serve alcuna gather finale, perché la matrice si compone
incrementalmente sul master. Vale la pena notare il contrasto con OpenMP: lì la
coda era un dettaglio interno al runtime, invisibile nel codice utente; qui la
stessa politica richiede un protocollo di messaggi esplicito, un master dedicato
e una condizione di terminazione da dimostrare. È il costo di complessità, prima
ancora che di tempo.

\paragraph{Strato driver.}
MPI richiede un sottile adattamento del driver condiviso, isolato dietro
\texttt{\#ifdef USE\_MPI} così da lasciare invariate le build seriale e OpenMP:
inizializzazione e chiusura di MPI, una barriera che allinea l'inizio della
regione cronometrata su tutti i rank, e l'I/O confinato al rank $0$, l'unico a
detenere la matrice assemblata. Il tempo di comunicazione raggiunge il CSV
attraverso l'accessore \texttt{kernel\_comm\_seconds()}, nullo per seriale e
OpenMP, senza che il kernel calcoli metriche.

\paragraph{Configurazione della misura: un solo nodo, per vincolo.}
Il job gira su un \textbf{unico nodo}, per un vincolo dell'ambiente e non per
scelta metodologica (sezione~\ref{sec:ambiente}): una richiesta con
\texttt{-{}-nodes=2} viene rifiutata da SLURM al momento della sottomissione. Il
paradigma nato per la memoria distribuita viene
dunque misurato, per forza di cose, su memoria condivisa.

L'allocazione è di $33$ task su un nodo, con i rank ancorati ai core e assegnati
in ordine, riempiendo un socket prima di passare all'altro
(\texttt{-{}-map-by core -{}-bind-to core}). È l'analogo MPI del
\texttt{OMP\_PROC\_BIND=close} usato per OpenMP: senza, il runtime potrebbe
distribuire i due soli rank di $p=2$ su socket diversi, caricando i punti a
basso $p$ di una penalità inter-socket estranea alla decomposizione. Soltanto a
$np=33$ un rank risiede necessariamente sul secondo socket. Si noti che questo posizionamento \emph{rank\,$\to$\,core}
è cosa distinta dalla decomposizione \emph{riga\,$\to$\,rank} che porta lo stesso
nome ``a blocchi''.

La conseguenza va dichiarata apertamente. Tutti i messaggi transitano per la
memoria condivisa del nodo e non per la rete, dunque la
\texttt{comm\_fraction} misurata è un \emph{limite inferiore ottimistico}
rispetto a quanto pagherebbe un'esecuzione realmente distribuita: quantifica
correttamente \emph{quanta} comunicazione ciascuno schema richiede, ma ne
sottostima il prezzo unitario. Il divario fra master--worker e schemi statici va
perciò letto come ottimistico verso il primo. In cambio si ottiene che $T_1$ e
$T_\parallel(p)$ sono garantiti sullo stesso hardware per costruzione, e che le
curve MPI e OpenMP diventano direttamente confrontabili nei valori assoluti, non
solo nelle tendenze: la loro differenza riflette in larga misura il solo
modello di programmazione.

\paragraph{Il lanciatore: \texttt{mpirun} e non \texttt{srun}.}
La versione di OpenMPI disponibile sul cluster non può essere lanciata
direttamente da \texttt{srun}: con il modulo \texttt{amd/gcc-8.5.0/openmpi-4.1.6}
ogni rank termina dentro \texttt{MPI\_Init} con l'errore
\begin{verbatim}
OPAL ERROR: Unreachable in file ext3x_client.c at line 111
\end{verbatim}
La causa è un disallineamento di versione nello strato di avvio: il client PMIx
di OpenMPI~$4.1.6$ è costruito per PMIx~$3.x$ e non dialoga con il server PMIx
v5 esposto da SLURM. I rank sono quindi lanciati con \texttt{mpirun}, che legge
l'allocazione dalle variabili d'ambiente di SLURM e avvia i processi con il PMIx
interno a OpenMPI. Il posizionamento dei rank resta invariato
(\texttt{-{}-map-by core -{}-bind-to core}), mentre la variabile
\texttt{MPI\_DECOMP} va inoltrata esplicitamente con \texttt{-x}, perché
\texttt{mpirun} propaga di default le sole variabili \texttt{OMPI\_*}. La
baseline seriale, che non usa MPI, continua a essere lanciata con \texttt{srun}.

Restano valide le scelte già motivate per OpenMP e per l'ambiente di misura
(sezione~\ref{sec:ambiente}): $T_1$ misurato con il binario seriale vero sullo
stesso nodo, compilazione dentro il job, kernel \emph{bare} come riferimento di
scaling.

% ============================================================
\subsubsection{Analisi dei risultati}
\label{sec:mpi-risultati}
% ============================================================

% Dati: job 34137 su gnode01, partizione ulow, 53 s di elapsed, COMPLETED 0:0.
% run_34137.csv (15 punti MPI + la propria baseline seriale) e merged.csv.
% T(1) = 0,720205 s, misurato nello stesso job.

Il confronto mette a sistema i tre schemi sullo stesso intervallo di unità di
calcolo, a griglia e $N_{\max}$ fissati. La tabella~\ref{tab:mpi-schemi}
raccoglie i $15$ punti misurati.\footnote{Job $34137$ su \texttt{gnode01},
    partizione \texttt{ulow}. Gli speedup sono riferiti alla baseline seriale
    misurata \emph{dentro lo stesso job}, $T_1 = 0{,}72021$~s; lo sweep OpenMP usa
    analogamente la propria, $0{,}72563$~s. Le due differiscono dello $0{,}75\%$, che
    è il fondo di risoluzione di qualunque confronto \emph{fra} paradigmi condotto
    sugli speedup: per questo il raffronto diretto con OpenMP, più avanti, è fatto
    sui tempi.}
\vspace{2,5mm}
\begin{table}[htbp]
    \centering
    \begin{tabular}{llrrrrrr}
        \hline
        Schema & $p$ & $S \le p/\lambda_p$ & $T_\parallel(p)$ [s] & $S(p)$ & $E(p)$ & $e(p)$ & comm.\ [\%] \\
        \hline
        seriale  & 1  & ---     & $0{,}72021$ & --- & --- & --- & --- \\
        \hline
        block    & 2  & $2{,}00$  & $0{,}36113$ & $1{,}994$  & $0{,}997$ & $0{,}0029$ & $0{,}13$ \\
        block    & 4  & $2{,}13$  & $0{,}33965$ & $2{,}120$  & $0{,}530$ & $0{,}2955$ & $0{,}08$ \\
        block    & 8  & $3{,}41$  & $0{,}21255$ & $3{,}388$  & $0{,}424$ & $0{,}1944$ & $0{,}14$ \\
        block    & 16 & $6{,}17$  & $0{,}11759$ & $6{,}125$  & $0{,}383$ & $0{,}1075$ & $0{,}31$ \\
        block    & 32 & $12{,}02$ & $0{,}06157$ & $11{,}697$ & $0{,}366$ & $0{,}0560$ & $1{,}20$ \\
        \hline
        cyclic   & 2  & $2{,}00$  & $0{,}36252$ & $1{,}987$  & $0{,}993$ & $0{,}0067$ & $0{,}18$ \\
        cyclic   & 4  & $4{,}00$  & $0{,}18126$ & $3{,}973$  & $0{,}993$ & $0{,}0022$ & $0{,}25$ \\
        cyclic   & 8  & $7{,}99$  & $0{,}09096$ & $7{,}918$  & $0{,}990$ & $0{,}0015$ & $0{,}24$ \\
        cyclic   & 16 & $15{,}97$ & $0{,}04930$ & $14{,}608$ & $0{,}913$ & $0{,}0064$ & $3{,}37$ \\
        cyclic   & 32 & $31{,}70$ & $0{,}02411$ & $29{,}871$ & $0{,}933$ & $0{,}0023$ & $3{,}02$ \\
        \hline
        dynamic  & 2  & $2{,}00$  & $0{,}36199$ & $1{,}990$  & $0{,}995$ & $0{,}0053$ & $0{,}43$ \\
        dynamic  & 4  & $4{,}00$  & $0{,}18154$ & $3{,}967$  & $0{,}992$ & $0{,}0028$ & $0{,}68$ \\
        dynamic  & 8  & $8{,}00$  & $0{,}09092$ & $7{,}921$  & $0{,}990$ & $0{,}0014$ & $0{,}86$ \\
        dynamic  & 16 & $16{,}00$ & $0{,}04622$ & $15{,}583$ & $0{,}974$ & $0{,}0018$ & $2{,}21$ \\
        dynamic  & 32 & $32{,}00$ & $0{,}02316$ & $\mathbf{31{,}102}$ & $0{,}972$ & $0{,}0009$ & $2{,}54$ \\
        \hline
    \end{tabular}
    \caption{MPI: metriche per schema di decomposizione e numero di unità di
    calcolo. Griglia $1024^2$, $N_{\max}=1000$, $33$ rank su un unico nodo. Per lo
    schema dinamico $p$ indica i worker; i processi allocati sono $p+1$. La colonna
        $S \le p/\lambda_p$ riporta il limite predetto a priori dal profilo di carico
        seriale (tabella~\ref{tab:serial_pred_lambdaP}) e non è una misura: è il termine
        di confronto contro cui va letto lo speedup osservato.}
    \label{tab:mpi-schemi}
\end{table}
\vspace{2,5mm}
\begin{table}[htbp]
    \centering
    \begin{tabular}{lrrrr}
        \hline
        $p$ & $T$ OpenMP \texttt{dynamic,1} [s] & $T$ MPI \texttt{dynamic} [s] &
        divario & comm.\ MPI [\%] \\
        \hline
        2  & $0{,}36079$ & $0{,}36199$ & $+0{,}33\%$ & $0{,}43$ \\
        4  & $0{,}18043$ & $0{,}18154$ & $+0{,}62\%$ & $0{,}68$ \\
        8  & $0{,}09028$ & $0{,}09092$ & $+0{,}71\%$ & $0{,}86$ \\
        16 & $0{,}04542$ & $0{,}04622$ & $+1{,}76\%$ & $2{,}21$ \\
        32 & $0{,}02342$ & $0{,}02316$ & $-1{,}11\%$ & $2{,}54$ \\
        \hline
    \end{tabular}
    \caption{Confronto dei due scheduler dinamici a parità di politica (concessione
    di una riga): stessa strategia su modelli di memoria diversi. Il confronto è sui
    \emph{tempi} e non sugli speedup, così da essere indipendente dalla baseline
    seriale usata dai due sweep. Divario positivo significa MPI più lento. Si ricordi
    che a $p$ unità di calcolo MPI occupa $p+1$ core, OpenMP $p$.}
    \label{tab:dynamic-confronto}
\end{table}

\paragraph{Il quadro d'insieme.}
A $32$ unità di calcolo, sullo stesso hardware e sullo stesso carico, lo speedup
va da $11{,}70$ (block) a $31{,}10$ (master--worker): un fattore $2{,}66$ fra la
regola peggiore e la migliore, che coincide entro lo $0{,}1\%$ con lo
sbilanciamento $\lambda_{32}$ previsto per i blocchi dal solo profilo seriale.
Il rapporto fra speedup misurato e limite previsto è compreso fra $0{,}972$ e
$0{,}997$ in tredici punti su quindici; le due eccezioni sono il cyclic a $p=16$
e $p=32$ ($0{,}915$ e $0{,}942$), discusse più avanti. Il confronto punto per
punto, esteso a OpenMP, è nella sezione~\ref{sec:validazione-previsioni}.

\paragraph{Il block è limitato dal proprio $\lambda_p$.}
Il dato più vistoso della tabella è il passaggio da $p=2$ a $p=4$: raddoppiando
i processi il tempo scende da $0{,}36113$~s a $0{,}33965$~s, cioè del $6\%$.
Non è un'anomalia di misura: è $\lambda_p$ che passa da $1{,}0000$ a $1{,}8798$
e si mangia per intero il raddoppio. Il punto $p=2$ merita una nota, perché è
ingannevole: la sua efficienza quasi perfetta ($0{,}997$) non dice che il block
scali bene a $p$ piccolo, ma che la viewport è simmetrica rispetto all'asse
reale, sicché le due metà dell'immagine contengono esattamente lo stesso numero
di iterazioni ($129\,827\,745$ ciascuna) e $\lambda_2 = 1$ per una proprietà
geometrica dell'insieme, non della decomposizione.

Al crescere di $p$ il $\lambda_p$ dei blocchi continua a crescere, fino al valore
intrinseco $\lambda = 2{,}76$ del caso limite $p = H$. L'efficienza tende perciò
a stabilizzarsi attorno a $1/\lambda_p$, pari a $0{,}376$ a $p=32$ contro lo
$0{,}366$ misurato, mentre lo speedup continua a crescere: già a $p=32$, però,
circa il $63\%$ del tempo di calcolo aggregato va perso in attesa.

\paragraph{Lettura di Karp--Flatt.}
Il block è il caso di scuola in cui la metrica segnala correttamente un problema
attribuendolo però alla causa sbagliata. Se il deficit fosse una frazione
seriale alla Amdahl, $e(p)$ sarebbe costante; qui invece decresce di un fattore
cinque, da $0{,}2955$ a $0{,}0560$, e segue la forma
$e(p) \approx (\lambda_p - 1)/(p - 1)$ già ricavata per lo \texttt{static} di
OpenMP (sezione~\ref{sec:openmp-risultati}), che realizza la stessa regola: i
valori predetti, $0{,}2933$, $0{,}1927$, $0{,}1062$ e $0{,}0536$ per
$p = 4, 8, 16, 32$, riproducono i misurati entro l'$1\%$ fino a $p=16$; a
$p=32$ lo scarto sale al $4{,}5\%$, nello stesso punto in cui lo speedup si
allontana dal limite. La ``frazione seriale
sperimentale'' non misura codice seriale, ma sbilanciamento.

\paragraph{Block e cyclic: bilanciamento a costo zero.}
Il primo confronto isola il puro effetto del bilanciamento: la comunicazione è
la stessa unica gather, dunque ogni differenza di tempo è attribuibile alla sola
distribuzione delle righe. Fino a $p=8$ il cyclic è indistinguibile dall'ideale
--- $E = 0{,}993$, $0{,}993$, $0{,}990$ a $2$, $4$ e $8$ processi --- mentre il
block, sulla stessa gather e sullo stesso carico, si ferma a $0{,}997$, $0{,}530$
e $0{,}424$. A $p=8$ il rapporto fra i due tempi è $2{,}34$, contro il $2{,}35$
predetto dal rapporto dei rispettivi $\lambda_p$: l'intero divario è
distribuzione delle righe. La \texttt{comm\_fraction} chiude l'argomento: $0{,}18\%$,
$0{,}25\%$, $0{,}24\%$ per il cyclic contro $0{,}13\%$, $0{,}08\%$, $0{,}14\%$
per il block, cioè fra $0{,}2$ e $0{,}7$~ms in entrambi i casi --- la stessa cosa
entro il rumore di misura. Non c'è stato alcun dialogo aggiuntivo: solo la
medesima gather finale. L'equivalenza vale però soltanto su questo intervallo di
$p$, e si rompe a $p=16$: se ne discute subito sotto.

\paragraph{Il master--worker e il costo del coordinamento.}
Il master--worker firma il miglior punto assoluto dell'intera campagna CPU:
$S = 31{,}10$ su $32$ unità di calcolo, $E = 0{,}972$, $89{,}7$~GFLOP/s contro i
$2{,}88$ della baseline seriale. Il residuo del $2{,}8\%$ è comunicazione, e la
colonna che lo misura è già in tabella. Confrontando riga per riga il deficit
dal limite $p/\lambda_p$ con la \texttt{comm\_fraction} si ottiene $0{,}5\%$
contro $0{,}43\%$ a $p=2$, poi $0{,}8$ contro $0{,}68$, $1{,}0$ contro $0{,}86$,
$2{,}6$ contro $2{,}21$ e $2{,}8\%$ contro $2{,}54\%$ a $p=32$: i due andamenti
coincidono entro pochi decimi di punto su tutti e cinque i punti. Lo stesso
numero si ricava per via indipendente dal protocollo: a $p=32$ un worker spende
circa $18~\mu$s di round-trip per riga contro $0{,}70$~ms di calcolo medio, cioè
il $2{,}6\%$ del proprio tempo. La coincidenza non è casuale ma strutturale:
poiché il bilanciamento è perfetto ($\lambda_p = 1$) e il tempo di un worker è
calcolo più attesa, il deficit dal limite \emph{è} la frazione di comunicazione,
e non resta un residuo da attribuire ad altro.

Vale la pena escludere esplicitamente la spiegazione alternativa, perché è quella
che il modello del makespan suggerirebbe. Il lavoro non è frazionabile a piacere:
l'unità indivisibile è la riga, e l'ultima riga consegnata a un worker potrebbe
in linea di principio essere la più pesante, lasciando gli altri fermi ---
l'\emph{effetto coda}. Su una consegna in ordine arbitrario il sovraccosto
sarebbe di al più $\lambda_{\text{riga}}\,p/H$, cioè circa l'$8{,}6\%$ a $p=32$.
Qui però le righe sono consegnate in ordine $0 \ldots H-1$, e la coda si svuota
perciò sulle \emph{ultime} righe della griglia, che sono di bordo e fra le più
leggere dell'immagine ($2\,281$ iterazioni contro una media di $253\,570$). Le
righe pesanti vengono distribuite a metà esecuzione, quando ogni altro worker ha
ancora lavoro da prelevare e nessuno resta inattivo. Simulando lo scheduling
greedy sul profilo $\{w_r\}$ misurato, con la consegna nell'ordine che il codice
usa davvero, il makespan supera la quota equa $\mathcal{W}/p$ dello $0{,}013\%$
a $p=32$: in questo problema e con questo ordine l'effetto coda è oltre due
ordini di grandezza più piccolo del deficit osservato, e non lo spiega.

Va ribadita la contabilità: a $32$ worker lo schema alloca $33$ core, uno dei
quali non calcola. L'efficienza per unità di calcolo è $0{,}972$, quella per core
allocato $31{,}10/33 = 0{,}942$. Entrambi i numeri sono legittimi, ma vanno
dichiarati insieme: a parità di core \emph{allocati}, e non di unità di calcolo,
il vantaggio del master--worker sul cyclic si assottiglia sensibilmente.

\paragraph{Cyclic e master--worker: quando il coordinamento adattivo comincia a valere.}
Fino a $p=8$ i due schemi sono indistinguibili: $0{,}09096$~s contro
$0{,}09092$~s, uno scarto dello $0{,}04\%$, un ordine di grandezza sotto la
dispersione delle ripetizioni; a $p=4$ il cyclic è addirittura più veloce dello
$0{,}15\%$. Su questo intervallo la conclusione forte regge per intero: il
self-scheduling non ripaga la propria complessità, perché una decomposizione
statica accorta bilancia già quanto basta. A $p=16$ e $p=32$, però, il
master--worker impiega il $6{,}3\%$ e il $4{,}0\%$ di tempo in meno del cyclic.

Il deficit del cyclic ad alto $p$ non è imputabile a $\lambda_p$, che è quasi
$1$. Va invece imputato, in parti misurabili, a due costi che il cyclic paga e
gli altri due schemi no.

Il primo è la gather, che ad alto $p$ non è più la stessa del block. In valore
assoluto, a $p=16$ il cyclic vi spende $1{,}66$~ms contro i $0{,}36$~ms del
block, un fattore $4{,}6$: quel divario di $1{,}3$~ms copre da solo circa il
$40\%$ dei $3{,}1$~ms che separano il cyclic dal master--worker. A $p=32$,
invece, le due gather tornano equivalenti ($0{,}73$ contro $0{,}74$~ms), il che
isola ulteriormente $p=16$ come il punto anomalo della campagna.

Il secondo è tutto ciò che il rank $0$ esegue serialmente attorno alla gather. Il
cyclic, oltre alla matrice finale, alloca un secondo buffer da $4$~MiB in cui
raccogliere le righe impacchettate, e alla fine ne copia $4$~MiB nella matrice
riga per riga (\repofile{mandelbrot/src/mandelbrot_mpi.cpp}): due azzeramenti e
una copia completa contro l'unico azzeramento del block, tutti di costo
indipendente da $p$. È precisamente la firma che la metrica di Karp--Flatt sa
riconoscere, e in questo schema la riconosce. La $e(p)$ del cyclic non decresce
come quella del block né cresce come quella del dinamico: resta \emph{costante}
attorno a $0{,}0015$--$0{,}0023$, esclusi i due punti rumorosi ($p=2$, dove pesa
lo scarto fra baseline, e $p=16$). Una $e(p)$ costante è la firma di Amdahl, e
$0{,}002$ su $T_1 = 720$~ms valgono circa $1{,}4$~ms di codice genuinamente
seriale, coerente con la stima del riordino e delle allocazioni. È l'unico punto
dell'intera campagna in cui la ``frazione seriale sperimentale'' misura davvero
codice seriale, e non sbilanciamento o coordinamento travestiti.

Resta che a $p=16$ i due termini bastano appena: $1{,}3$~ms di gather in eccesso
più circa $1{,}4$~ms di costi seriali coprono $2{,}7$ dei $3{,}1$~ms di divario,
e quel punto è anche quello con la dispersione più alta. La verifica diretta ---
cronometrare il riordino, oppure eliminarlo ricevendo in posizione con un tipo
derivato \texttt{MPI\_Type\_vector} --- rientra fra gli sviluppi del lavoro,
insieme alla rimisura del punto.

\paragraph{La comunicazione è piccola, ma è ciò che resta.}
In termini assoluti la comunicazione è modesta in tutta la campagna: fra
$0{,}22$ e $1{,}66$~ms su tempi che vanno da $23$~ms a $361$~ms, mai oltre il
$3{,}4\%$ di $T_\parallel(p)$. Il master--worker scambia $2H = 2\,048$ messaggi
per ripetizione contro l'unica gather collettiva degli schemi statici, e
nonostante ciò resta sotto il $2{,}6\%$: alla grana di una riga ($0{,}70$~ms di
calcolo contro $18~\mu$s di protocollo) il messaggio non impedisce allo schema
di avvicinarsi all'ideale. Sarebbe però sbagliato concluderne che sia
trascurabile: una volta che il bilanciamento è perfetto non resta altro, ed è
esattamente questa quota a separare il master--worker dallo speedup lineare.

Tre cautele accompagnano però questi numeri. Primo, \texttt{comm\_fraction} non è
la stessa grandezza nei tre schemi: negli statici misura la sola
\texttt{MPI\_Gatherv}, nel dinamico include anche l'attesa del comando successivo; confrontare il
$2{,}54\%$ del dinamico con l'$1{,}20\%$ del block come se fossero omogenei è
scorretto. Secondo, per il block quasi tutto il tempo perso non è in
\texttt{comm}: a $p=32$ lo schema ha $1{,}20\%$ di comunicazione e $63\%$ di
inattività parcheggiata nella barriera, conteggiata in $T_\parallel(p)$ ma non
in \texttt{comm\_fraction}; letta da sola, la colonna suggerirebbe che il block
è efficiente. Terzo, questi valori sono su memoria condivisa: con la rete di
mezzo il master--worker, \emph{latency-bound} su $2\,048$ messaggi piccoli,
degraderebbe assai più del cyclic, \emph{bandwidth-bound} su una singola
collettiva --- al punto che la classifica potrebbe invertirsi. È la conclusione
più esposta al limite del nodo singolo.

\paragraph{Scheduler dinamici a confronto: il costo del modello di memoria.}
Il terzo confronto, in tabella~\ref{tab:dynamic-confronto}, mette a fronte i due
scheduler dinamici sulla stessa politica: a concessione unitaria da entrambi i
lati, il loro divario quantifica il costo del solo cambio di modello di memoria.
Fino a $p=16$ MPI paga, e paga poco: dallo $0{,}33\%$ all'$1{,}76\%$. È il
prezzo del protocollo rispetto a un prelievo da contatore condiviso, e la sua
crescita con $p$ è quella attesa da un coordinatore centralizzato.

A $p=32$ il segno si inverte: MPI è più veloce dell'$1{,}11\%$ ($0{,}02316$~s
contro $0{,}02342$~s). Da un lato, \texttt{dynamic,1} di OpenMP mostra proprio a
$32$ thread il suo unico scostamento apprezzabile dall'ideale ($E$ da $0{,}999$
a $0{,}968$): $1\,024$ prelievi da un unico contatore condiviso fra $32$ thread
cominciano a costare contesa. Dall'altro, una possibile spiegazione è che il master
MPI serializzi le richieste su un core proprio, senza contendere cache line ai
worker. Alla scala
massima misurata, il costo della coda in memoria condivisa supera quello del
passaggio di messaggi. Resta la contabilità: quel vantaggio si ottiene occupando
un $33^{\circ}$ core, e a parità di core allocati l'ordine si ristabilisce a
favore di OpenMP di circa il $2\%$.

\paragraph{Correttezza e qualità della misura.}
Tutte e sedici le righe del file di run condividono il medesimo checksum
($9203734293765563287$), uguale a quello della baseline seriale: le tre
decomposizioni riproducono la matrice degli escape time bit per bit, e nessuno
schema ha alterato il risultato. Un caso resta però scoperto, e va dichiarato:
tutti i $p$ misurati ($2 \ldots 32$) dividono $H = 1024$, e nello schema dinamico
la divisibilità è comunque irrilevante perché le righe sono consegnate una alla
volta. Il ramo che assegna una riga in più ai primi \texttt{rem} rank non è
dunque mai esercitato; coprirlo richiederebbe un punto a $p$ non divisore. Le sei baseline seriali a $1024^2$ e $N_{\max}=1000$ registrate nel
dataset, una per ciascun job, valgono fra $0{,}71990$ e $0{,}72563$~s e
concordano entro lo $0{,}8\%$: è una verifica indipendente della riproducibilità
del nodo. La dispersione fra ripetizioni, misurata come scarto della mediana dal
minimo, resta sotto il $2\%$ nella maggior parte dei punti, ma sale al $4{,}4\%$
per il master--worker a $p=32$ e al $5{,}1\%$ e $6{,}6\%$ per il cyclic a $p=16$
e $p=32$, fra i punti più brevi e quindi più esposti alla contesa su
\texttt{gnode01}. L'efficienza non monotona del cyclic ($0{,}913$ a $p=16$
contro $0{,}933$ a $p=32$) non ha spiegazione strutturale ed è il punto da
rimisurare prima di costruirvi sopra un'affermazione.